\Article{28}
Da die Formeln für $X$, $Y$, $Z$ zusammen 71 Glieder enthalten, ist die unmittelbare Rechnung nach diesen Formeln recht mühsam. Sie für eine große Zahl von Orten zu wiederholen, wäre umso abschreckender, als man kaum hoffen könnte, vor möglichen Rechenfehlern geschützt zu
\expandafter\def\csname HistoricalNoteText19\endcsname{Herr Dr. \textsc{Goldschmidt} hat einen Teil dieser Hilfstafel berechnet.}
\SourceLeaf{157}sein, ohne dieselbe Rechnung zweimal auszuführen. Auch würde man wenig gewinnen, wenn man alle Glieder wegließe, deren Koeffizienten weniger als eine Einheit oder sogar weniger als 10 Einheiten betragen, da immer noch 65 Glieder übrig blieben. Da aber der Wert der ganzen Arbeit ungewiss bliebe, wenn man sie nicht anhand einer beträchtlichen Zahl tatsächlicher Beobachtungen prüfte, habe ich die Mühe auf mich genommen, eine Hilfstafel zu berechnen\HistoricalNote{19}. Ihr Gebrauch verkürzt und erleichtert die Arbeit in hohem Maße und verbessert dadurch den Schutz vor Rechenfehlern wesentlich. Ihr Aufbau beruht darauf, dass die Werte der Komponenten in die folgende Form gebracht werden
\[
\begin{aligned}
X={}&a^0+a'\cos(\lambda+A')+a''\cos(2\lambda+A'')\\
&+a'''\cos(3\lambda+A''')+a''''\cos(4\lambda+A'''')\\
Y={}&b'\cos(\lambda+B')+b''\cos(2\lambda+B'')\\
&+b'''\cos(3\lambda+B''')+b''''\cos(4\lambda+B'''')\\
Z={}&c^0+c'\cos(\lambda+C')+c''\cos(2\lambda+C'')\\
&+c'''\cos(3\lambda+C''')+c''''\cos(4\lambda+C'''')
\end{aligned}
\]
Die erste Tafel enthält die von $\lambda$ unabhängigen Teile von $X$ und $Z$. In den vier folgenden findet man die Werte der Hilfswinkel $A'$, $A''$ usw. und der Logarithmen von $a'$, $a''$ usw., jeweils für die einzelnen Grade der Breite $\varphi=90^0-u$. [Die Tafel ist dem vorliegenden Abdruck am Ende dieser Abhandlung in einer wegen des anderen Formats etwas geänderten Anordnung beigefügt.]

Als Beispiel sei hier die Rechnung für Göttingen angeführt.

Für die Breite $+51^0\,32'$ findet man in den Tafeln
\[
\begin{array}{r|r|r}
a^0=+500{,}8&&c^0=+1465{,}2\\
\log a'=2{,}28980&\log b'=2{,}18900&\log c'=2{,}20204\\
\log a''=1{,}79403&\log b''=2{,}03220&\log c''=2{,}12777\\
\log a'''=1{,}32522&\log b'''=1{,}46845&\log c'''=1{,}43199\\
\log a''''=0{,}59391&\log b''''=0{,}70016&\log c''''=0{,}59091\\
A'=249^0\,30'&B'=358^0\,24'&C'=105^0\,44'\\
A''=311\quad45&B''=64\quad50&C''=165\quad15\\
A'''=234\quad10&B'''=318\quad13&C'''=42\quad22\\
A''''=142\quad26&B''''=232\quad26&C''''=322\quad26
\end{array}
\]
\SourceLeaf{158}
und daraus mit der Länge $\lambda=9^0\,56\frac12'$ die Beiträge zu
\[
\begin{array}{r|r|r}
X&Y&Z\\\hline
+500{,}8&&+1465{,}2\\
-35{,}71&+152{,}89&-68{,}99\\
+54{,}76&+9{,}92&-133{,}67\\
-2{,}21&+28{,}77&+8{,}27\\
-3{,}92&+0{,}19&+3{,}90\\\hline
X=+513{,}72&Y=+191{,}77&Z=+1274{,}71
\end{array}
\]
Die weitere Rechnung ergibt dann
\[
\begin{aligned}
\delta&=+20^0\,28'\quad\log\omega=2{,}73907\\
i&=+66\quad43\\
\psi&=1387{,}6\quad\text{oder in der gewöhnlichen Einheit}\\
\psi&=1{,}3876
\end{aligned}
\]

